\section{Fixed-fibre Newton geometry}\label{sec:fixednewton}

For a fixed direction \(v\), let
\[
K_v(\e,t)=\sum_{p,q\ge0}A_{p,q}(v)\e^pt^q.
\]
Its Newton polyhedron is the convex hull of
\(\{(p,q):A_{p,q}(v)\ne0\}+\R_{\ge0}^2\).  Actual coefficients, rather than
candidate support, define this set.

\begin{theorem}[Fixed-fibre Newton polyhedra]\label{thm:fibrenewton}
For the three branches of \cref{thm:trichotomy}, the Newton polyhedron is,
respectively,
\[
(1,3)+\R_{\ge0}^2,\qquad
(2,4)+\R_{\ge0}^2,\qquad
\varnothing.
\]
Every positive weight exposes the unique displayed vertex in the non-dark
branches.
\end{theorem}
\begin{proof}
The trichotomy supplies the vertex and its non-zero coefficient.  The endpoint
bound shows that every other surviving exponent lies in the corresponding
upper orthant.  A strictly positive linear functional increases along every
non-zero vector in that orthant, hence exposes the vertex uniquely.
\end{proof}

\noindent\textbf{Formal status.}
The definition uses the actual non-zero coefficient support from M6 and the
Mathlib real convex hull.  The empty-germ convention, the general one-vertex
orthant lemma, the three fibre shapes and positive-weight uniqueness are
\evidence{KERNEL-CHECKED} as \path{M3A.newtonPolyhedron_empty},
\path{M3A.newton_eq_orthant_of_vertex},
\path{M3A.fixedFibre_newton} and
\path{M3A.fixedFibre_positiveWeight_faces}.

\begin{figure}[ht]
\centering\begin{tikzpicture}[scale=.78,>=Latex]
\begin{scope}[xshift=-4.2cm]
\draw[->] (0,0)--(4.3,0) node[right] {$p$}; \draw[->] (0,0)--(0,4.3) node[above] {$q$};
\fill[blue!12] (1,3) rectangle (4,4); \fill[blue!12] (1,3) rectangle (4,3.25);
\draw[blue!60!black,thick] (1,3)--(4,3); \draw[blue!60!black,thick] (1,3)--(1,4);
\fill (1,3) circle (2.3pt); \node[below right,font=\small] at (1,3) {$(1,3)$};
\node at (2.1,-.7) {linear stratum};
\end{scope}
\begin{scope}[xshift=1cm]
\draw[->] (0,0)--(4.3,0) node[right] {$p$}; \draw[->] (0,0)--(0,4.3) node[above] {$q$};
\fill[orange!15] (2,4) rectangle (4,4.15); \draw[orange!80!black,thick] (2,4)--(4,4); \draw[orange!80!black,thick] (2,4)--(2,4.15);
\fill (2,4) circle (2.3pt); \node[below right,font=\small] at (2,4) {$(2,4)$};
\node at (2.1,-.7) {quadratic stratum};
\end{scope}
\end{tikzpicture}

\caption{Fixed-fibre Newton polyhedra.  The vertex moves from \((1,3)\) to
\((2,4)\) on the quadratic stratum; the exact-dark fibre has empty support.}
\label{fig:fixednewton}
\end{figure}

The change of variables \(x=z^{-1}\) in \cref{cor:transformed} is not cosmetic.
The resolvent expansion is a Laurent series at \(z=\infty\); multiplying by
\(x^{-1}\) aligns the coefficient of \(x^k\) with moment depth \(k\).  Omitting
that shift displaces every Newton point by one and corrupts the comparison with
the amplitude.

Classical Newton-polyhedron methods organise dominant exponents and exposed
faces \parencite{Kouchnirenko1976,Varchenko1976}.  Here the fibre theorem is
especially simple because there is one lower vertex.  The subtlety begins when
directions move or coordinates are pulled back, allowing several parent
exponents to collide.

\medskip
\noindent\textbf{Formalisation boundary.}
This completes the M0--M8 end-to-end Lean spine.  Beginning with
\cref{sec:pb1}, the paper moves to conventionally proved extensions supported
by exact symbolic or combinatorial checks.  No theorem in the remaining
sections is represented as part of the current Lean package unless explicitly
stated otherwise.
